% Seção V — Identificação e Descrição do Problema
\section{Identificação e Descrição do Problema}
\label{sec:problema}

\subsection{Contexto e Delimitação}
\label{sec:contexto}

[Descrever o ambiente e contexto onde o problema ocorre.
Delimitar o escopo do trabalho: o que está dentro e o que está fora do escopo.
Justificar as delimitações adotadas.]

\subsection{Problema Central}
\label{sec:problema-central}

[Formular uma pergunta de pesquisa clara ou uma declaração objetiva do problema.
Exemplo: ``Como é possível [ação] de forma [característica] para [público-alvo]?'']

\subsection{Requisitos do Sistema}
\label{sec:requisitos}

% Requisitos funcionais — o que o sistema deve fazer
\subsubsection{Requisitos Funcionais}
\label{sec:rf}

\begin{itemize}
  \item RF01: [Descrição do requisito funcional 1];
  \item RF02: [Descrição do requisito funcional 2];
  \item RF03: [Descrição do requisito funcional 3].
\end{itemize}

% Requisitos não funcionais — qualidades e restrições do sistema
\subsubsection{Requisitos Não Funcionais}
\label{sec:rnf}

\begin{itemize}
  \item RNF01: [Desempenho --- ex.: resposta em menos de 200\,ms];
  \item RNF02: [Usabilidade --- ex.: interface intuitiva para usuários não técnicos];
  \item RNF03: [Portabilidade --- ex.: executar em Linux e Windows sem recompilação].
\end{itemize}
